\section{Introducción}
\label{sec:introduccion}
\subsection{Contexto y motivación}
El proyecto se enmarca en la asignatura Biometría Informática del Máster Universitario en Visión Artificial de la Universidad Rey Juan Carlos. Su objeto de estudio es la evaluación de calidad de imágenes faciales y el papel de los modelos de visión-lenguaje en esta tarea.

El análisis distinguirá la calidad perceptual de la utilidad biométrica. La primera se refiere a la apariencia visual de la imagen; la segunda se estudiará en relación con su contribución al rendimiento de un sistema de reconocimiento facial. Se investigará si una explicación en lenguaje natural permite identificar factores de captura relevantes y ofrecer recomendaciones útiles.
\pendiente{Sustentar la motivación con fuentes originales sobre FIQA y su relación con los errores de reconocimiento.}

\subsection{Pregunta de investigación y objetivos}
\begin{quote}
¿Pueden los VLMs evaluar y explicar la calidad de una imagen facial de manera que sus resultados se correspondan con su utilidad real para el reconocimiento biométrico y no únicamente con su calidad perceptual?
\end{quote}
El objetivo general es revisar y comparar los enfoques de FIQA, prestando especial atención a las aportaciones y limitaciones de los VLMs. Se plantean los siguientes objetivos específicos:
\begin{enumerate}
\item Organizar los fundamentos de calidad de imagen, FIQA y modelos de visión-lenguaje.
\item Comparar métodos especializados y multimodales mediante evidencia publicada.
\item Examinar por separado utilidad biométrica, consistencia y exactitud de las explicaciones.
\item Identificar limitaciones, condiciones de comparación y líneas de investigación futuras.
\end{enumerate}

\subsection{Alcance y organización}
El trabajo será principalmente teórico y de revisión comparativa, con una extensión final aproximada de 15 páginas en formato de artículo científico. La ejecución de experimentos propios queda condicionada al alcance que se acuerde con el profesor y a los recursos disponibles. La presentación asociada, de unas 20 diapositivas, se preparará como un entregable independiente.

La sección~\ref{sec:arte} organizará el estado del arte; la sección~\ref{sec:evaluacion} recogerá datasets, protocolos y comparaciones; la sección~\ref{sec:conclusiones} sintetizará los hallazgos y las cuestiones abiertas. El foco será la calidad biométrica facial.
